% The Survivor's Number, Group 6
% Author: Princilia DACI
\documentclass[12pt,a4paper]{article}
\usepackage[latin1]{inputenc}
\usepackage{amsmath}
\usepackage{amsfonts}
\usepackage{amssymb}
\usepackage{graphicx}
\usepackage{multirow,multicol}
\usepackage{pifont,hyperref,lastpage,fancyhdr,movie15,float}


\usepackage[left=3cm,right=3cm,top=4cm,bottom=3cm]{geometry}

%%%%%%%%%%%%%%%%%%%%%%%%%%%%%%%%%%%%%%%%%%%%%%%%%%%%%%%%%%%%%%%%%%%%%%%%%
%%%%%%%%%%%%%%%%%%%%%%%%%% Footer and Header %%%%%%%%%%%%%%%%%%%%%%%%%%%%
\pagestyle{fancy}                                                     %%%
\fancyhf{} 		                                                      %%%
\lfoot{\tiny\textsf{\includegraphics[width=1.8cm]{AIMSSenegalLogo}    %%%
Po.Box~1418 Mbour-Thies,~phone~{(+221) 33 956 7693},~                 %%%
\url{http://www.aims-senegal.org}}}                                   %%%
\rfoot{\bfseries Page \thepage~of \pageref{LastPage}}                 %%%
\renewcommand{\footrulewidth}{1.pt}\renewcommand{\headrulewidth}{0pt} %%%
%%%%%%%%%%%%%%%%%%%%%%%%%%%%%%%%%%%%%%%%%%%%%%%%%%%%%%%%%%%%%%%%%%%%%%%%%
%%%%%%%%%%%%%%%%%%%%%%%%%%%%%%%%%%%%%%%%%%%%%%%%%%%%%%%%%%%%%%%%%%%%%%%%%

%%%%%%%%%%%%%%%%%%%%%%%%%%%%%%%%%%%%%%%%%%%%%%%%%%%%%%%%%%%%%%%%%%%%%%%%%
%%%%%%% Fill here information about the present assignment %%%%%%%%%%%%%%
\newcommand{\code}{Group 6}                                       %%%
\newcommand{\deadline}{01 Nov 2026}                                  %%%
\newcommand{\assignment}{Group project (Problem 0.6) ON Mathematical Problem Solving}			  %%% 
\newcommand{\lecturer}{Lecturer: Dr Audace A. V. Dossou-Olory}                     %%%
%%%%%%%%%%%%%%%%%%%%%%%%%%%%%%%%%%%%%%%%%%%%%%%%%%%%%%%%%%%%%%%%%%%%%%%%%
%%%%%%%%%%%%%%%%%%%%%%%%%%%%%%%%%%%%%%%%%%%%%%%%%%%%%%%%%%%%%%%%%%%%%%%%%

%%%%%%%%%%%%%%%%%%%%%%%%%%%%%%%%%%%%%%%%%%%%%%%%%%%%%%%%%%%%%%%%%%%%%%%%%
%%%%%%%%%%%%%%%%%%%%%%  Title at the first page  %%%%%%%%%%%%%%%%%%%%%%%%
\title{\vspace*{-4cm}\begin{minipage}{\textwidth}                     %%%
\begin{center}                                                        %%%
\begin{tabular}{|c|c|c|}                                              %%%
\hline\multicolumn{3}{|c|}{\bf\scriptsize\MakeUppercase\assignment}\\ %%%
\hline{\small Student's Code}&                                        %%%
\multirow{3}{7cm}{\includegraphics[width=7.5cm,height=2.3cm]{AIMSSenegalLogo}} %%%
& {\small Deadline}\\                                                 %%%
\cline{1-1}\cline{3-3}{\small\bf\code}&&{\small\bf\deadline} \\       %%%
\cline{1-1}\cline{3-3}{\small\today} &&{\small2026-2027}\\            %%%
\hline\multicolumn{3}{|r|}{\scriptsize\lecturer}\\\hline              %%%
\end{tabular}                                                         %%%
\end{center}                                                          %%%
\end{minipage}\hfill\date{}\vspace*{-1cm}}                            %%%
%%%%%%%%%%%%%%%%%%%%%%%%%%%%%%%%%%%%%%%%%%%%%%%%%%%%%%%%%%%%%%%%%%%%%%%%%
%%%%%%%%%%%%%%%%%%%%%%%%%%%%%%%%%%%%%%%%%%%%%%%%%%%%%%%%%%%%%%%%%%%%%%%%%

\newcommand{\K}{\mathbb{K}}
\newcommand{\R}{\mathbb{R}}
\newcommand{\C}{\mathbb{C}}

\newtheorem{theo}{Theorem}
\newtheorem{defi}{Definition}
\newenvironment{proof}[1][Proof.]{\textbf{#1~}}{\ \rule{0.5em}{0.5em}}

%%%%%%%%%%%%%%%%%%%%%%%%%%%%%%%%%%%%%%%%%%%%%%%%%%%%%%%%%%%%%%%%%%%%%%%%%
%%%%%%%%%%%%%%%%%%%%% Added for our report (Group 6) %%%%%%%%%%%%%%%%%%%%
\hypersetup{hidelinks,
  pdfauthor={Princilia DACI},
  pdftitle={The Survivor's Number},
  pdfsubject={Mathematical Problem Solving: Projects, Group 6},
  pdfcreator={Princilia DACI},
  pdfproducer={Princilia DACI}}
\usepackage{tikz}
\usepackage{caption}
\usepackage[below]{placeins}
\usepackage{needspace}
\usepackage{xcolor}
\definecolor{halfgray}{gray}{0.55}
\definecolor{ipython_frame}{RGB}{207, 207, 207}
\definecolor{ipython_bg}{RGB}{247, 247, 247}
\definecolor{ipython_red}{RGB}{186, 33, 33}
\definecolor{ipython_green}{RGB}{0, 128, 0}
\definecolor{ipython_cyan}{RGB}{64, 128, 128}
\definecolor{ipython_purple}{RGB}{170, 34, 255}

\usepackage{listings}
\lstset{
    breaklines=true,
    %
    extendedchars=true,
    literate=
    {á}{{\'a}}1 {é}{{\'e}}1 {í}{{\'i}}1 {ó}{{\'o}}1 {ú}{{\'u}}1
    {Á}{{\'A}}1 {É}{{\'E}}1 {Í}{{\'I}}1 {Ó}{{\'O}}1 {Ú}{{\'U}}1
    {à}{{\`a}}1 {è}{{\`e}}1 {ì}{{\`i}}1 {ò}{{\`o}}1 {ù}{{\`u}}1
    {À}{{\`A}}1 {È}{{\'E}}1 {Ì}{{\`I}}1 {Ò}{{\`O}}1 {Ù}{{\`U}}1
    {ä}{{\"a}}1 {ë}{{\"e}}1 {ï}{{\"i}}1 {ö}{{\"o}}1 {ü}{{\"u}}1
    {Ä}{{\"A}}1 {Ë}{{\"E}}1 {Ï}{{\"I}}1 {Ö}{{\"O}}1 {Ü}{{\"U}}1
    {â}{{\^a}}1 {ê}{{\^e}}1 {î}{{\^i}}1 {ô}{{\^o}}1 {û}{{\^u}}1
    {Â}{{\^A}}1 {Ê}{{\^E}}1 {Î}{{\^I}}1 {Ô}{{\^O}}1 {Û}{{\^U}}1
    {œ}{{\oe}}1 {Œ}{{\OE}}1 {æ}{{\ae}}1 {Æ}{{\AE}}1 {ß}{{\ss}}1
    {ç}{{\c c}}1 {Ç}{{\c C}}1 {ø}{{\o}}1 {å}{{\r a}}1 {Å}{{\r A}}1
    {€}{{\EUR}}1 {£}{{\pounds}}1
}

%%
%% Python definition (c) 1998 Michael Weber
%% Additional definitions (2013) Alexis Dimitriadis
%% modified by me (should not have empty lines)
%%
\lstdefinelanguage{iPython}{
    morekeywords={access,and,break,class,continue,def,del,elif,else,except,exec,finally,for,from,global,if,import,in,is,lambda,not,or,pass,print,raise,return,try,while},%
    %
    % Built-ins
    morekeywords=[2]{abs,all,any,basestring,bin,bool,bytearray,callable,chr,classmethod,cmp,compile,complex,delattr,dict,dir,divmod,enumerate,eval,execfile,file,filter,float,format,frozenset,getattr,globals,hasattr,hash,help,hex,id,input,int,isinstance,issubclass,iter,len,list,locals,long,map,max,memoryview,min,next,object,oct,open,ord,pow,property,range,raw_input,reduce,reload,repr,reversed,round,set,setattr,slice,sorted,staticmethod,str,sum,super,tuple,type,unichr,unicode,vars,xrange,zip,apply,buffer,coerce,intern},%
    %
    sensitive=true,%
    morecomment=[l]\#,%
    morestring=[b]',%
    morestring=[b]",%
    %
    morestring=[s]{'''}{'''},% used for documentation text (mulitiline strings)
    morestring=[s]{"""}{"""},% added by Philipp Matthias Hahn
    %
    morestring=[s]{r'}{'},% `raw' strings
    morestring=[s]{r"}{"},%
    morestring=[s]{r'''}{'''},%
    morestring=[s]{r"""}{"""},%
    morestring=[s]{u'}{'},% unicode strings
    morestring=[s]{u"}{"},%
    morestring=[s]{u'''}{'''},%
    morestring=[s]{u"""}{"""},%
    %
    % {replace}{replacement}{lenght of replace}
    % *{-}{-}{1} will not replace in comments and so on
    literate=
    *{+}{{{\color{ipython_purple}+}}}1
    {-}{{{\color{ipython_purple}-}}}1
    {*}{{{\color{ipython_purple}$^\ast$}}}1
    {/}{{{\color{ipython_purple}/}}}1
    {^}{{{\color{ipython_purple}\^{}}}}1
    {?}{{{\color{ipython_purple}?}}}1
    {!}{{{\color{ipython_purple}!}}}1
    {\%}{{{\color{ipython_purple}\%}}}1
    {<}{{{\color{ipython_purple}<}}}1
    {>}{{{\color{ipython_purple}>}}}1
    {|}{{{\color{ipython_purple}|}}}1
    {\&}{{{\color{ipython_purple}\&}}}1
    {~}{{{\color{ipython_purple}~}}}1
    %
    {==}{{{\color{ipython_purple}==}}}2
    {<=}{{{\color{ipython_purple}<=}}}2
    {>=}{{{\color{ipython_purple}>=}}}2
    %
    {+=}{{{+=}}}2
    {-=}{{{-=}}}2
    {*=}{{{$^\ast$=}}}2
    {/=}{{{/=}}}2,
    %
    literate=
    {á}{{\'a}}1 {é}{{\'e}}1 {í}{{\'i}}1 {ó}{{\'o}}1 {ú}{{\'u}}1
    {Á}{{\'A}}1 {É}{{\'E}}1 {Í}{{\'I}}1 {Ó}{{\'O}}1 {Ú}{{\'U}}1
    {à}{{\`a}}1 {è}{{\`e}}1 {ì}{{\`i}}1 {ò}{{\`o}}1 {ù}{{\`u}}1
    {À}{{\`A}}1 {È}{{\'E}}1 {Ì}{{\`I}}1 {Ò}{{\`O}}1 {Ù}{{\`U}}1
    {ä}{{\"a}}1 {ë}{{\"e}}1 {ï}{{\"i}}1 {ö}{{\"o}}1 {ü}{{\"u}}1
    {Ä}{{\"A}}1 {Ë}{{\"E}}1 {Ï}{{\"I}}1 {Ö}{{\"O}}1 {Ü}{{\"U}}1
    {â}{{\^a}}1 {ê}{{\^e}}1 {î}{{\^i}}1 {ô}{{\^o}}1 {û}{{\^u}}1
    {Â}{{\^A}}1 {Ê}{{\^E}}1 {Î}{{\^I}}1 {Ô}{{\^O}}1 {Û}{{\^U}}1
    {œ}{{\oe}}1 {Œ}{{\OE}}1 {æ}{{\ae}}1 {Æ}{{\AE}}1 {ß}{{\ss}}1
    {ç}{{\c c}}1 {Ç}{{\c C}}1 {ø}{{\o}}1 {å}{{\r a}}1 {Å}{{\r A}}1
    {€}{{\EUR}}1 {£}{{\pounds}}1,
    %
%   identifierstyle=\color{red}\ttfamily,
    commentstyle=\color{ipython_cyan}\ttfamily,
    stringstyle=\color{ipython_red}\ttfamily,
    keepspaces=true,
    showspaces=false,
    showstringspaces=false,
    %
    rulecolor=\color{ipython_frame},
    frame=single,
    frameround={t}{t}{t}{t},
    framexleftmargin=6mm,
    numbers=left,
    numberstyle=\tiny\color{halfgray},
    %
    %
    backgroundcolor=\color{ipython_bg},
    %   extendedchars=true,
    basicstyle=\scriptsize\ttfamily,
    keywordstyle=\color{ipython_green}\ttfamily,
}
\lstset{columns=fullflexible,keepspaces=true,basicstyle=\footnotesize\ttfamily}
\newtheorem{obs}{Observation}
\newtheorem{lemma}{Lemma}
\newtheorem{coro}{Corollary}
% no single line left alone at the top or bottom of a page
\clubpenalty=10000
\widowpenalty=10000

% people on a circle (same colours as in our presentation)
\providecolor{aimscolour}{rgb}{0.54,0.11,0.01}
\newcommand{\nodesize}{5mm}
\newcommand{\nodefont}{\scriptsize}
\tikzset{
  person/.style={circle,draw,thick,fill=white,minimum size=\nodesize,inner sep=0pt,font=\nodefont},
  gone/.style={circle,draw=gray!60,fill=white,text=gray!60,minimum size=\nodesize,inner sep=0pt,font=\nodefont}
}
% cross on an eliminated person; the optional argument is the colour
\newcommand{\cross}[2][aimscolour]{\draw[#1,very thick] (#2.north west) -- (#2.south east) (#2.north east) -- (#2.south west);}
% n, radius, people eliminated in earlier rounds, people eliminated in this round, survivor (0 if not shown)
\newcommand{\gamecircle}[5]{%
  \begin{tikzpicture}[baseline=(current bounding box.center)]
    \foreach \p in {1,...,#1}
      \node[person] (p\p) at ({90-360/#1*(\p-1)}:#2) {\p};
    \foreach \e in {#3} { \node[gone] (b\e) at (p\e) {\e}; \cross[gray!40]{b\e} }
    \foreach \e in {#4} { \node[gone] (o\e) at (p\e) {\e}; \cross{o\e} }
    \ifnum#5>0 \draw[aimscolour,very thick] (p#5) circle ({0.5*\nodesize+0.8mm}); \fi
  \end{tikzpicture}}
%%%%%%%%%%%%%%%%%%%%%%%%%%%%%%%%%%%%%%%%%%%%%%%%%%%%%%%%%%%%%%%%%%%%%%%%%

\begin{document}
\maketitle\thispagestyle{fancy}

% first page: title and group
\vspace*{2.5cm}
\begin{center}
{\LARGE\bf The Survivor's Number}\\[12pt]
{\large Problem 0.6 of the Mathematical Problem Solving projects}\\[2.5cm]
{\Large\bf Group 6}\\[12pt]
\renewcommand{\arraystretch}{1.3}
\begin{tabular}{ll}
Daci Yimete Princilia Doleresse & Dlamini Sandziso Siyabonga\\
Fadhul Abdouroihamane & Kambale Maruba Exauc\'e\\
Marie Hortense Lucie Senamie Saizonou & Oumar Sow\\
Razan Salah Mohammed Salih Mansour & Samba Diallo
\end{tabular}\\[2.5cm]
Supervised by Dr Audace A. V. Dossou-Olory\\[6pt]
African Institute for Mathematical Sciences (AIMS), Senegal\\[6pt]
November 2026
\end{center}
\newpage

% second page: abstract
\begin{abstract}
People numbered $1$ to $n$ stand in a circle. Starting with person $1$, every
second person is eliminated until only one is left, and $J(n)$ denotes the
number of this survivor. We first played the game by hand for
$n=1,\dots,16$. The results suggested three rules: $J(2k)=2J(k)-1$,
$J(2k+1)=2J(k)+1$, and the fact that every $n$ can be written in exactly one
way as $n=2^m+\ell$ with $0\leq\ell<2^m$. We proved them with one lemma: the
game only depends on the places of the people around the circle, not on their
numbers, so after one turn of the circle it becomes a smaller copy of itself.
We then proved the formula
\[
J(2^m+\ell)=2\ell+1
\]
in two ways, by induction on $m$ and by induction on $n$.
We generalised it to a game that starts with any person $s$: the survivor is
the person $2\ell$ places after $s$. In base $2$, the formula becomes a simple
rule, since $J(n)$ is obtained by moving the first digit of $n$ to the end.
This gives at once, for example, $J(10^9)=926\,258\,177$. Finally, we explain
the Python code we wrote to simulate the game and to check every result.
\end{abstract}
\newpage

% third page: contents
\setcounter{tocdepth}{2}
\tableofcontents
\newpage

\section{Introduction}

This report presents the work of Group~6 on Problem~0.6 of the project sheet
of the course Mathematical Problem Solving \cite{projects2026,course2026},
entitled \emph{Survivor's number}.

\subsection{The problem}

The statement given in the project sheet \cite{projects2026} is the
following.

\begin{center}\fbox{\begin{minipage}{0.94\textwidth}
We start with $n$ people numbered $1$ to $n$ around a circle, and we
eliminate every second remaining person until only one survives. For example,
for $n=10$, the elimination order is $2,4,6,8,10$, then $3,7,1,9$, so $5$
survives. Determine the survivor's number, $J(n)$. For example, $J(10)=5$.
The case $n=0$ makes no sense. First, check that
\[
J(1)=1,\quad J(2)=1,\quad J(3)=3,\quad J(4)=1,\quad J(5)=3,\quad J(6)=5 .
\]
\end{minipage}}\end{center}

We make the rules precise.

\begin{defi}[The elimination game]\label{def:game}
Let $n\geq1$. We place $n$ people, numbered $1,2,\dots,n$, around a circle,
and after person $n$ comes person $1$ again. We go around the circle, always in
the same direction, and apply the following rules:
\begin{itemize}
  \item we skip one person, then we eliminate the next one, then we skip the
        next one, and so on;
  \item an eliminated person leaves the circle, so only the people still
        present are counted;
  \item the game stops when only one person is left.
\end{itemize}
The number of the last person is the \emph{survivor's number}.
\end{defi}

This elimination game is a classical problem, treated for example in
\cite[Chapter~1]{concrete}.

In the example of the statement, with $n=10$, the game skips $1$ and
eliminates $2$, skips $3$ and eliminates $4$, and so on until $10$. It then
comes back to the beginning of the circle and continues with the people still
present, until only person $5$ is left.

\subsection{Why we need a formula}

For ten people, the game can be played on paper in a minute. With $1000$
people it needs $999$ eliminations, and after the first turn the people left
are no longer consecutive numbers, so it becomes hard to follow who comes
next. So the question is: \emph{can we find $J(n)$ without playing the whole
game?} Our objectives are to
\begin{enumerate}
  \item find a formula that gives $J(n)$ directly for every $n\geq1$, and prove
        it;
  \item find the survivor when the game starts with any person;
  \item find the survivor quickly for very large circles, and check our
        results with a computer program.
\end{enumerate}

\subsection{Who starts the game}\label{sec:assumption}

In the statement, the first person eliminated is person $2$. So person $1$ is
the first person skipped: the game starts with person $1$. Until the end of
Section~\ref{sec:formula} we keep this assumption, and $J(n)$ always means the
survivor of this game. This is not a restriction: a game that starts with
another person is the same game turned around the circle. Once $J(n)$ is
known, we use this remark in Section~\ref{sec:general} to treat any starting
person.

\subsection{Method and plan of the report}

We followed the method of the course \cite{course2026}: play with small
cases, look for a pattern, make a precise guess, and prove it so that it holds
for every $n$ and not only for the cases we tested. The report follows the
same order.
\begin{itemize}
  \item \textbf{Section 2} plays the game for $n=1,\dots,16$ and collects the
        results in a table.
  \item \textbf{Section 3} proves a key lemma, which says that the game only
        depends on the places of the people around the circle. With it, we
        prove how the survivor for $2k$ and $2k+1$ people depends on the
        survivor for $k$ people, and we show that every $n$ is written in one
        way as a power of $2$ plus a remainder.
  \item \textbf{Section 4} proves the formula for $J(n)$ in two ways, by
        induction on $m$ and by induction on $n$, gives its consequences, and
        generalises it to a game that starts with any person.
  \item \textbf{Section 5} writes the formula in base $2$, which gives the
        survivor at once for very large circles.
  \item \textbf{Section 6} explains the Python code we wrote to play the game
        and to check our results.
  \item \textbf{Section 7} concludes.
\end{itemize}

\clearpage
\section{Playing the Game by Hand}

Before looking for a formula, we play the game by hand for every $n$ from $1$
to $16$, starting with person~$1$ as in Section~\ref{sec:assumption}. The results are
collected in Table~\ref{tab:results} at the end of this section.

\subsection{Conventions}\label{sec:conventions}

A \emph{round} is one turn of the circle. It ends when the game comes back to
the beginning of the circle, and the next round goes on with the people still
present. When $n$ is odd, the first round skips person $n$ and eliminates
person $1$, who comes right after him: we count this elimination in the first
round.

Up to $n=7$, and for $n=16$, we draw the circle at the end of each round and
write under it the people eliminated during that round, in order. From $n=8$
to $n=15$ we draw only the final circle, with the rounds written under it.
The symbols are the following:
\begin{center}
\begin{tikzpicture}[baseline=-0.6ex]
  \node[person] at (0,0) {4};
  \node[anchor=west,font=\footnotesize] at (0.35,0) {still in the circle};
  \node[gone] (a) at (3.4,0) {2}; \cross{a}
  \node[anchor=west,font=\footnotesize] at (3.75,0) {eliminated};
  \node[gone] (b) at (6.2,0) {2}; \cross[gray!40]{b}
  \node[anchor=west,font=\footnotesize] at (6.55,0) {eliminated in an earlier round};
  \node[person] (c) at (11.9,0) {5}; \draw[aimscolour,very thick] (c) circle ({0.5*\nodesize+0.8mm});
  \node[anchor=west,font=\footnotesize] at (12.3,0) {survivor};
\end{tikzpicture}
\end{center}

\subsection{The sixteen tests}

\subsubsection{For $n=1,2,3$}
\par\vspace{6pt}\noindent
\begin{minipage}[t]{0.32\linewidth}\centering
\textbf{$n=1$}\\[3pt]
{\setlength{\tabcolsep}{3pt}\begin{tabular}{p{2.40cm}}
\centering\gamecircle{1}{0}{}{}{1}\tabularnewline[2pt]
\centering\scriptsize No elimination\\ $J(1)=1$\tabularnewline
\end{tabular}}
\end{minipage}\hfill
\begin{minipage}[t]{0.32\linewidth}\centering
\textbf{$n=2$}\\[3pt]
{\setlength{\tabcolsep}{3pt}\begin{tabular}{p{2.40cm}}
\centering\gamecircle{2}{0.45}{}{2}{1}\tabularnewline[2pt]
\centering\scriptsize Round 1: 2\\ $J(2)=1$\tabularnewline
\end{tabular}}
\end{minipage}\hfill
\begin{minipage}[t]{0.32\linewidth}\centering
\textbf{$n=3$}\\[3pt]
{\setlength{\tabcolsep}{3pt}\begin{tabular}{p{2.40cm}}
\centering\gamecircle{3}{0.60}{}{2,1}{3}\tabularnewline[2pt]
\centering\scriptsize Round 1: 2, 1\\ $J(3)=3$\tabularnewline
\end{tabular}}
\end{minipage}\par\vspace{6pt}

\subsubsection{For $n=4$ to $7$}
\par\vspace{6pt}\noindent
\begin{minipage}[t]{0.48\linewidth}\centering
\textbf{$n=4$}\\[3pt]
{\setlength{\tabcolsep}{3pt}\begin{tabular}{p{2.40cm}p{2.40cm}}
\centering\gamecircle{4}{0.60}{}{2,4}{0} & \centering\gamecircle{4}{0.60}{2,4}{3}{1}\tabularnewline[2pt]
\centering\scriptsize Round 1: 2, 4 & \centering\scriptsize Round 2: 3\\ $J(4)=1$\tabularnewline
\end{tabular}}
\end{minipage}\hfill
\begin{minipage}[t]{0.48\linewidth}\centering
\textbf{$n=5$}\\[3pt]
{\setlength{\tabcolsep}{3pt}\begin{tabular}{p{2.40cm}p{2.40cm}}
\centering\gamecircle{5}{0.85}{}{2,4,1}{0} & \centering\gamecircle{5}{0.85}{2,4,1}{5}{3}\tabularnewline[2pt]
\centering\scriptsize Round 1: 2, 4, 1 & \centering\scriptsize Round 2: 5\\ $J(5)=3$\tabularnewline
\end{tabular}}
\end{minipage}\par\vspace{6pt}

\par\vspace{6pt}\noindent
\begin{minipage}[t]{0.48\linewidth}\centering
\textbf{$n=6$}\\[3pt]
{\setlength{\tabcolsep}{3pt}\begin{tabular}{p{2.40cm}p{2.40cm}}
\centering\gamecircle{6}{0.85}{}{2,4,6}{0} & \centering\gamecircle{6}{0.85}{2,4,6}{3,1}{5}\tabularnewline[2pt]
\centering\scriptsize Round 1: 2, 4, 6 & \centering\scriptsize Round 2: 3, 1\\ $J(6)=5$\tabularnewline
\end{tabular}}
\end{minipage}\hfill
\begin{minipage}[t]{0.48\linewidth}\centering
\textbf{$n=7$}\\[3pt]
{\setlength{\tabcolsep}{3pt}\begin{tabular}{p{2.40cm}p{2.40cm}}
\centering\gamecircle{7}{0.85}{}{2,4,6,1}{0} & \centering\gamecircle{7}{0.85}{2,4,6,1}{5,3}{7}\tabularnewline[2pt]
\centering\scriptsize Round 1: 2, 4, 6, 1 & \centering\scriptsize Round 2: 5, 3\\ $J(7)=7$\tabularnewline
\end{tabular}}
\end{minipage}\par\vspace{6pt}

\subsubsection{For $n=8$ to $15$}
\par\vspace{6pt}\noindent
\begin{minipage}[t]{0.32\linewidth}\centering
\textbf{$n=8$}\\[3pt]
{\gamecircle{8}{1.00}{}{2,4,6,8,3,7,5}{1}}\\[3pt]
{\scriptsize Round 1: 2, 4, 6, 8\\ Round 2: 3, 7\\ Round 3: 5\\ $J(8)=1$}
\end{minipage}\hfill
\begin{minipage}[t]{0.32\linewidth}\centering
\textbf{$n=9$}\\[3pt]
{\gamecircle{9}{1.00}{}{2,4,6,8,1,5,9,7}{3}}\\[3pt]
{\scriptsize Round 1: 2, 4, 6, 8, 1\\ Round 2: 5, 9\\ Round 3: 7\\ $J(9)=3$}
\end{minipage}\hfill
\begin{minipage}[t]{0.32\linewidth}\centering
\textbf{$n=10$}\\[3pt]
{\gamecircle{10}{1.00}{}{2,4,6,8,10,3,7,1,9}{5}}\\[3pt]
{\scriptsize Round 1: 2, 4, 6, 8, 10\\ Round 2: 3, 7, 1\\ Round 3: 9\\ $J(10)=5$}
\end{minipage}\par\vspace{6pt}

\par\vspace{6pt}\noindent
\begin{minipage}[t]{0.32\linewidth}\centering
\textbf{$n=11$}\\[3pt]
{\gamecircle{11}{1.00}{}{2,4,6,8,10,1,5,9,3,11}{7}}\\[3pt]
{\scriptsize Round 1: 2, 4, 6, 8, 10, 1\\ Round 2: 5, 9, 3\\ Round 3: 11\\ $J(11)=7$}
\end{minipage}\hfill
\begin{minipage}[t]{0.32\linewidth}\centering
\textbf{$n=12$}\\[3pt]
{\gamecircle{12}{1.00}{}{2,4,6,8,10,12,3,7,11,5,1}{9}}\\[3pt]
{\scriptsize Round 1: 2, 4, 6, 8, 10, 12\\ Round 2: 3, 7, 11\\ Round 3: 5, 1\\ $J(12)=9$}
\end{minipage}\hfill
\begin{minipage}[t]{0.32\linewidth}\centering
\textbf{$n=13$}\\[3pt]
{\renewcommand{\nodesize}{4.6mm}\gamecircle{13}{1.15}{}{2,4,6,8,10,12,1,5,9,13,7,3}{11}}\\[3pt]
{\scriptsize Round 1: 2, 4, 6, 8, 10, 12, 1\\ Round 2: 5, 9, 13\\ Round 3: 7, 3\\ $J(13)=11$}
\end{minipage}\par\vspace{6pt}

\par\vspace{6pt}\noindent
\hfill
\begin{minipage}[t]{0.32\linewidth}\centering
\textbf{$n=14$}\\[3pt]
{\renewcommand{\nodesize}{4.6mm}\gamecircle{14}{1.15}{}{2,4,6,8,10,12,14,3,7,11,1,9,5}{13}}\\[3pt]
{\scriptsize Round 1: 2, 4, 6, 8, 10, 12, 14\\ Round 2: 3, 7, 11, 1\\ Round 3: 9, 5\\ $J(14)=13$}
\end{minipage}\hfill
\begin{minipage}[t]{0.32\linewidth}\centering
\textbf{$n=15$}\\[3pt]
{\renewcommand{\nodesize}{4.6mm}\gamecircle{15}{1.15}{}{2,4,6,8,10,12,14,1,5,9,13,3,11,7}{15}}\\[3pt]
{\scriptsize Round 1: 2, 4, 6, 8, 10, 12, 14, 1\\ Round 2: 5, 9, 13, 3\\ Round 3: 11, 7\\ $J(15)=15$}
\end{minipage}\hfill\null
\par\vspace{6pt}

\subsubsection{For $n=16$}
\par\vspace{6pt}\noindent
\begin{minipage}[t]{1\linewidth}\centering
\textbf{$n=16$}\\[3pt]
{\renewcommand{\nodesize}{4.3mm}\renewcommand{\nodefont}{\tiny}\setlength{\tabcolsep}{3pt}\begin{tabular}{p{3.00cm}p{3.00cm}p{3.00cm}p{3.00cm}}
\centering\gamecircle{16}{1.15}{}{2,4,6,8,10,12,14,16}{0} & \centering\gamecircle{16}{1.15}{2,4,6,8,10,12,14,16}{3,7,11,15}{0} & \centering\gamecircle{16}{1.15}{2,4,6,8,10,12,14,16,3,7,11,15}{5,13}{0} & \centering\gamecircle{16}{1.15}{2,4,6,8,10,12,14,16,3,7,11,15,5,13}{9}{1}\tabularnewline[2pt]
\centering\scriptsize Round 1: 2, 4, 6, 8, 10, 12, 14, 16 & \centering\scriptsize Round 2: 3, 7, 11, 15 & \centering\scriptsize Round 3: 5, 13 & \centering\scriptsize Round 4: 9\\ $J(16)=1$\tabularnewline
\end{tabular}}
\end{minipage}\par\vspace{6pt}

\subsection{Summary of the tests}

We write each $n$ as the largest power of $2$ that is not bigger than $n$,
plus a remainder: $n=2^m+\ell$ with $0\leq\ell<2^m$. Table~\ref{tab:results}
puts this decomposition next to the survivor found in the tests.

\begin{center}
\renewcommand{\arraystretch}{1.0}
\begin{tabular}{|c|c|c|c|c|c|}
\hline
$n$ & $n=2^m+\ell$ & $m$ & $\ell$ & $J(n)$ & $2\ell+1$\\ \hline
$1$ & $2^{0}+0$ & $0$ & $0$ & $1$ & $1$\\
\hline
$2$ & $2^{1}+0$ & $1$ & $0$ & $1$ & $1$\\
$3$ & $2^{1}+1$ & $1$ & $1$ & $3$ & $3$\\
\hline
$4$ & $2^{2}+0$ & $2$ & $0$ & $1$ & $1$\\
$5$ & $2^{2}+1$ & $2$ & $1$ & $3$ & $3$\\
$6$ & $2^{2}+2$ & $2$ & $2$ & $5$ & $5$\\
$7$ & $2^{2}+3$ & $2$ & $3$ & $7$ & $7$\\
\hline
$8$ & $2^{3}+0$ & $3$ & $0$ & $1$ & $1$\\
$9$ & $2^{3}+1$ & $3$ & $1$ & $3$ & $3$\\
$10$ & $2^{3}+2$ & $3$ & $2$ & $5$ & $5$\\
$11$ & $2^{3}+3$ & $3$ & $3$ & $7$ & $7$\\
$12$ & $2^{3}+4$ & $3$ & $4$ & $9$ & $9$\\
$13$ & $2^{3}+5$ & $3$ & $5$ & $11$ & $11$\\
$14$ & $2^{3}+6$ & $3$ & $6$ & $13$ & $13$\\
$15$ & $2^{3}+7$ & $3$ & $7$ & $15$ & $15$\\
\hline
$16$ & $2^{4}+0$ & $4$ & $0$ & $1$ & $1$\\
\hline
\end{tabular}
\captionof{table}{The tests for $n=1,\dots,16$. The horizontal lines separate
the blocks $2^m\leq n<2^{m+1}$.}
\label{tab:results}
\end{center}

Three things stand out in Table~\ref{tab:results}.
\begin{itemize}
  \item All the survivors are odd.
  \item The survivor is person $1$ when $n=1,2,4,8,16$, that is when $n$ is a
        power of $2$.
  \item Inside each block between two lines, the survivor goes up by $2$ each
        time a person is added. In every line, $J(n)$ is equal to $2\ell+1$.
\end{itemize}
These are only remarks on sixteen cases. We prove them in Sections~3
and~4.

\clearpage
\section{What the Table Shows, and Why}

Table~\ref{tab:results} suggests several rules. In this section we state them
one by one and prove each of them, so that they hold for every number of
people and not only for the sixteen circles we played. We first prove a lemma
that we use in all the proofs: the game only depends on the places of the
people, not on their numbers. The first two rules then link a circle of $2k$
or $2k+1$ people to a circle of $k$ people. The third one is about the way
each $n$ is written in the table.

\subsection{The key lemma: only the places matter}\label{sec:tool}

At any moment of the game, the people still present stand around the circle
in a fixed order. We write them as a list, starting with the person who is
skipped next.

\begin{defi}\label{def:list}
A \emph{list} $(a_1,a_2,\dots,a_k)$ is the sequence of the $k$ people still
present, in the order of the circle, where $a_1$ is the next person to be
skipped. If $k\geq2$, one \emph{step} of the game skips $a_1$ and eliminates
$a_2$. The next person to be skipped is then $a_3$, and $a_1$ is now the last
person before we come back to $a_3$:
\[
(a_1,a_2,a_3,\dots,a_k)\longrightarrow(a_3,a_4,\dots,a_k,a_1).
\]
For $k=2$, the new list is $(a_1)$. When $k=1$ the game is over and $a_1$ is
the survivor.
\end{defi}

Each step removes one person, so a game with $k$ people ends after $k-1$
steps. The game of Definition~\ref{def:game} starts from the list
$(1,2,\dots,n)$. For example, with $n=5$:
\[
(1,2,3,4,5)\to(3,4,5,1)\to(5,1,3)\to(3,5)\to(3).
\]
The people eliminated are $2,4,1,5$, in this order, and the survivor is $3$,
as in Table~\ref{tab:results}.

The step rule only uses the places in the list, never the numbers of the
people. This is what the next lemma says.

\begin{lemma}[Only the places matter]\label{lem:places}
Let $(a_1,\dots,a_k)$ be a list of $k\geq1$ people. The survivor of the game
that starts from this list is $a_{J(k)}$, the person at place $J(k)$.
\end{lemma}

\begin{proof}
We prove by induction on $k\geq1$ the statement
\[
L(k):\quad\text{for every list $(a_1,\dots,a_k)$ of $k$ people, the survivor
is $a_{J(k)}$.}
\]
\emph{Base case.} For $k=1$, the only person $a_1$ is the survivor, and
$J(1)=1$. So $L(1)$ is true.

\emph{Induction step.} Let $k\geq2$ and assume $L(k-1)$. Take a list
$(a_1,\dots,a_k)$. One step gives the list of $k-1$ people
\[
(b_1,\dots,b_{k-1})=(a_3,a_4,\dots,a_k,a_1),
\]
that is, $b_j=a_{f(j)}$ with
\[
f(j)=\begin{cases} j+2 & \text{if } 1\leq j\leq k-2,\\ 1 & \text{if } j=k-1.
\end{cases}
\]
The game goes on from this new list, so by $L(k-1)$ its survivor is
$b_{J(k-1)}=a_{f(J(k-1))}$. This is true for every list. For the list
$(1,2,\dots,k)$, which is the usual game with $k$ people, the survivor is
$J(k)$ by definition, so $J(k)=f\big(J(k-1)\big)$. Going back to any list, the
survivor is therefore $a_{f(J(k-1))}=a_{J(k)}$, and $L(k)$ is true.

By induction, $L(k)$ is true for every $k\geq1$.
\end{proof}

The proof also gives a relation between $J(k)$ and $J(k-1)$. Since
$J(k)=f\big(J(k-1)\big)$,
\begin{equation}\label{eq:onestep}
J(k)=\begin{cases} J(k-1)+2 & \text{if } J(k-1)\leq k-2,\\
1 & \text{if } J(k-1)=k-1 .\end{cases}
\end{equation}
In words: adding one person moves the survivor two places, except when the
survivor was the last person, in which case it goes back to person $1$. This
is what we saw in Table~\ref{tab:results}. Relation~\eqref{eq:onestep} gives
$J(n)$ for every $n$, but it needs $n-1$ steps. The two relations of the next
subsections are much faster, because they divide $n$ by $2$.

\subsection{Even circles: \texorpdfstring{$J(2k)=2J(k)-1$}{J(2k) = 2J(k) - 1}}\label{sec:even}

\subsubsection{What we observe}

We put side by side the survivor for $k$ people and the survivor for $2k$
people.
\begin{center}
\renewcommand{\arraystretch}{1.2}
\begin{tabular}{|c|c|c|c|c|c|c|c|c|}
\hline
$k$ & 1 & 2 & 3 & 4 & 5 & 6 & 7 & 8\\ \hline
$J(k)$ & 1 & 1 & 3 & 1 & 3 & 5 & 7 & 1\\ \hline
$J(2k)$ & 1 & 1 & 5 & 1 & 5 & 9 & 13 & 1\\ \hline
$2J(k)-1$ & 1 & 1 & 5 & 1 & 5 & 9 & 13 & 1\\ \hline
\end{tabular}
\end{center}
The last two lines are equal in all eight cases.

\begin{obs}\label{obs:even}
For every $k\geq1$, \ $J(2k)=2J(k)-1$.
\end{obs}

\subsubsection{Proof}

Let $k\geq1$, and let $n=2k$ people stand around the circle. The idea is to
play the first round completely, and to see that what is left is a game with
$k$ people.

\textbf{Step 1: the first round eliminates exactly the even numbers.} We show
by induction on $j$, for $j=1,\dots,k$, the statement
\[
E(j):\quad\text{the first $j$ eliminations skip $1,3,\dots,2j-1$ and
eliminate $2,4,\dots,2j$.}
\]
\emph{Base case.} The game starts by skipping person $1$ and eliminating
person $2$, so $E(1)$ is true.

\emph{Induction step.} Let $1\leq j<k$ and assume $E(j)$. At this point the
people eliminated are $2,4,\dots,2j$, and the last one is $2j$. The next
person is $2j+1$. The people
$2j+1$ and $2j+2$ are still present, and $2j+2\leq2k$. So the game skips
$2j+1$ and eliminates $2j+2$, and $E(j+1)$ is true.

By $E(k)$, the $k$-th person eliminated is $2k$, the last person of the circle. The
next person is person $1$, who is still present: we are back at the beginning
of the circle and the first round is over. It has eliminated
$2,4,\dots,2k$, and the $k$ people left are, in the order of the circle,
\[
1,\ 3,\ 5,\ \dots,\ 2k-1 .
\]
The game goes on by skipping person $1$.

\textbf{Step 2: the people left form a game with $k$ people.} We write the
people left as a list $(a_1,a_2,\dots,a_k)$, in the order of the circle,
starting with the person who is skipped next:
\[
a_i=2i-1\qquad(i=1,\dots,k),
\]
so $a_1=1$, $a_2=3$, \dots, $a_k=2k-1$ (Figure~\ref{fig:even}). The rest of the
game is the game of Definition~\ref{def:game} played on $a_1,\dots,a_k$ instead
of $1,\dots,k$: the circle has $k$ people and it starts by skipping the first
one.

\textbf{Step 3: the survivor is $a_{J(k)}$.} By Lemma~\ref{lem:places},
applied to the list $(a_1,\dots,a_k)$, the survivor is the person at place
$J(k)$ of this list. So
\[
J(2k)=a_{J(k)}=2J(k)-1 .
\]
This proves Observation~\ref{obs:even} for every $k\geq1$. \hfill$\square$

\begin{figure}[!ht]
\centering
\renewcommand{\nodesize}{7mm}\renewcommand{\nodefont}{\footnotesize}\footnotesize
\begin{tikzpicture}
  \draw[gray,dotted,very thick] (-35:2.1) arc (-35:-235:2.1);
  \node[person] (a1) at (90:2.1) {1};
  \node[gone] (a2) at (55:2.1) {}; \cross{a2} \node[fill=white,inner sep=0.4pt] at (a2) {2};
  \node[person] (a3) at (20:2.1) {3};
  \node[gone] (a4) at (-15:2.1) {}; \cross{a4} \node[fill=white,inner sep=0.4pt] at (a4) {4};
  \node[gone] (a6) at (125:2.1) {}; \cross{a6} \node[font=\scriptsize,fill=white,inner sep=0.4pt] at (a6) {$2k$};
  \node[person,font=\scriptsize,minimum size=8.5mm] at (160:2.1) {$2k{-}1$};
  \node[font=\scriptsize] at (90:2.75) {$a_1$};
  \node[font=\scriptsize] at (20:2.75) {$a_2$};
  \node[font=\scriptsize] at (160:2.8) {$a_k$};
\end{tikzpicture}
\caption{The circle of $2k$ people after the first round: $a_i=2i-1$. The
dotted arc stands for the people not drawn.}
\label{fig:even}
\end{figure}

\textbf{Example.} For $n=10$, the first round leaves $(1,3,5,7,9)$. Since
$J(5)=3$, the survivor is the third person of this list, $a_3=5$. So
$J(10)=2\times3-1=5$, as in Table~\ref{tab:results}.

\subsection{Odd circles: \texorpdfstring{$J(2k+1)=2J(k)+1$}{J(2k+1) = 2J(k) + 1}}

\subsubsection{What we observe}

We now compare the survivor for $k$ people and for $2k+1$ people.
\begin{center}
\renewcommand{\arraystretch}{1.2}
\begin{tabular}{|c|c|c|c|c|c|c|c|}
\hline
$k$ & 1 & 2 & 3 & 4 & 5 & 6 & 7\\ \hline
$J(k)$ & 1 & 1 & 3 & 1 & 3 & 5 & 7\\ \hline
$J(2k+1)$ & 3 & 3 & 7 & 3 & 7 & 11 & 15\\ \hline
$2J(k)+1$ & 3 & 3 & 7 & 3 & 7 & 11 & 15\\ \hline
\end{tabular}
\end{center}
Again, the last two lines are equal in all seven cases.

\begin{obs}\label{obs:odd}
For every $k\geq1$, \ $J(2k+1)=2J(k)+1$.
\end{obs}

\subsubsection{Proof}

Let $k\geq1$, and let $n=2k+1$ people stand around the circle. We follow the
same three steps as for $2k$ people. Only the end of the first round changes.

\textbf{Step 1: the first round.} The induction of the previous proof works
without change: the first $k$ eliminations skip $1,3,\dots,2k-1$ and
eliminate $2,4,\dots,2k$. Indeed, for $j<k$ we still have $2j+2\leq2k<2k+1$.

After the $k$-th elimination (person $2k$), one person remains before the end
of the circle: person $2k+1$. He is skipped. The next person is person $1$,
since after person $2k+1$ we come back to the beginning of the circle, and
person $1$ is eliminated. As said in Section~\ref{sec:conventions}, this elimination belongs to
the first round. So the first round eliminates $2,4,\dots,2k$ and then $1$,
and the $k$ people left are
\[
3,\ 5,\ 7,\ \dots,\ 2k+1 .
\]
The game goes on with the first person after $1$ who is still present.
Person $2$ is already eliminated, so this is person $3$, and he is the next one
skipped.

\textbf{Step 2: the people left form a game with $k$ people.} This time the
list starts with person $3$:
\[
a_i=2i+1\qquad(i=1,\dots,k),
\]
so $a_1=3$, $a_2=5$, \dots, $a_k=2k+1$ (Figure~\ref{fig:odd}). The rest of the
game is again the game of Definition~\ref{def:game} with $k$ people, played on
$a_1,\dots,a_k$.

\textbf{Step 3: the survivor.} By Lemma~\ref{lem:places}, the survivor is
the person at place $J(k)$ of the list, so
\[
J(2k+1)=a_{J(k)}=2J(k)+1 .
\]
This proves Observation~\ref{obs:odd} for every $k\geq1$. \hfill$\square$

\begin{figure}[!ht]
\centering
\renewcommand{\nodesize}{7mm}\renewcommand{\nodefont}{\footnotesize}\footnotesize
\begin{tikzpicture}
  \draw[gray,dotted,very thick] (-35:2.1) arc (-35:-235:2.1);
  \node[gone] (b1) at (90:2.1) {}; \cross{b1} \node[fill=white,inner sep=0.4pt] at (b1) {1};
  \node[gone] (b2) at (55:2.1) {}; \cross{b2} \node[fill=white,inner sep=0.4pt] at (b2) {2};
  \node[person] (b3) at (20:2.1) {3};
  \node[gone] (b4) at (-15:2.1) {}; \cross{b4} \node[fill=white,inner sep=0.4pt] at (b4) {4};
  \node[gone] (b5) at (160:2.1) {}; \cross{b5} \node[font=\scriptsize,fill=white,inner sep=0.4pt] at (b5) {$2k$};
  \node[person,font=\scriptsize,minimum size=8.5mm] at (125:2.1) {$2k{+}1$};
  \node[font=\scriptsize] at (20:2.75) {$a_1$};
  \node[font=\scriptsize] at (125:2.8) {$a_k$};
\end{tikzpicture}
\caption{The circle of $2k+1$ people after the first round: $a_i=2i+1$.
Person $1$ is eliminated at the end of the round, so the list starts with
person $3$.}
\label{fig:odd}
\end{figure}

\textbf{Example.} For $n=11$, the first round eliminates $2,4,6,8,10$ and then
$1$, and leaves $(3,5,7,9,11)$. Since $J(5)=3$, the survivor is $a_3=7$. So
$J(11)=2\times3+1=7$, as in Table~\ref{tab:results}.

\FloatBarrier
\subsubsection{Computing \texorpdfstring{$J(n)$}{J(n)} by halving}\label{sub:halving}

With $J(1)=1$, Observations~\ref{obs:even} and~\ref{obs:odd} give $J(n)$ for
every $n$. We divide $n$ by $2$ (rounding down) until we reach $1$, then we go
back up with the two relations. For $n=100$:
\[
100\to50\to25\to12\to6\to3\to1,
\]
and
\begin{align*}
J(3)&=2J(1)+1=3, & J(6)&=2J(3)-1=5, & J(12)&=2J(6)-1=9,\\
J(25)&=2J(12)+1=19, & J(50)&=2J(25)-1=37, & J(100)&=2J(50)-1=73.
\end{align*}
For $n=1000$, the halvings are
$1000\to500\to250\to125\to62\to31\to15\to7\to3\to1$, and going back up gives
\begin{align*}
&J(3)=3, && J(7)=7, && J(15)=15, && J(31)=31, && J(62)=61,\\
&J(125)=123, && J(250)=245, && J(500)=489, && J(1000)=977 .
\end{align*}
About $\log_2 n$ steps replace the $n-1$ eliminations of the game: $6$ steps
for $n=100$ and $9$ steps for $n=1000$. Still, this does not tell us what $J(n)$
looks like. For that we go back to the way the numbers are written in
Table~\ref{tab:results}.

\subsection{Every \texorpdfstring{$n$}{n} is \texorpdfstring{$2^m+\ell$}{2\^{}m+l} in exactly one way}

\subsubsection{What we observe}

In Table~\ref{tab:results}, each $n$ is written as $2^m+\ell$, where $2^m$ is
the largest power of $2$ that is not bigger than $n$:
\[
5=2^2+1,\qquad 10=2^3+2,\qquad 15=2^3+7,\qquad 16=2^4+0 .
\]
In every line the remainder $\ell$ is smaller than $2^m$, and there is only
one such way to write $n$. For instance, $10=2^2+6$ is not allowed, because
$6\geq2^2$.

\begin{obs}\label{obs:decomp}
Every integer $n\geq1$ can be written in exactly one way as
\[
n=2^m+\ell,\qquad m\geq0,\quad 0\leq\ell<2^m .
\]
\end{obs}

\subsubsection{Proof by induction}

\textbf{Existence, by strong induction on $n$.} For $n=1$ we have
$1=2^0+0$, and $0<2^0$.

Let $n\geq2$, and assume that every integer $k$ with $1\leq k<n$ can be
written $k=2^p+r$ with $0\leq r<2^p$. We write $n=2k$ if $n$ is even, or
$n=2k+1$ if $n$ is odd. In both cases $1\leq k<n$, so $k=2^p+r$ with
$0\leq r<2^p$.
\begin{itemize}
  \item If $n=2k$, then $n=2^{p+1}+2r$, and $0\leq2r<2\cdot2^p=2^{p+1}$.
  \item If $n=2k+1$, then $n=2^{p+1}+(2r+1)$. Since $r\leq2^p-1$, we get
        $2r+1\leq2^{p+1}-1<2^{p+1}$.
\end{itemize}
In both cases $n=2^m+\ell$ with $m=p+1$ and $0\leq\ell<2^m$.

\textbf{Uniqueness.} The condition $0\leq\ell<2^m$, with $\ell=n-2^m$, means
\begin{equation}\label{eq:block}
2^m\leq n<2^{m+1},
\end{equation}
since $\ell\geq0$ means $2^m\leq n$, and $\ell<2^m$ means
$n<2^m+2^m=2^{m+1}$. Assume that $n=2^m+\ell=2^{m'}+\ell'$ are two such
writings, with $m<m'$. Then $m+1\leq m'$, and \eqref{eq:block} for $m$ and for
$m'$ gives
\[
n<2^{m+1}\leq2^{m'}\leq n,
\]
which is impossible. In the same way $m'<m$ is impossible. So $m=m'$, and
then $\ell=n-2^m=\ell'$. \hfill$\square$

The proof also shows how to find $m$ and $\ell$: by \eqref{eq:block}, $2^m$ is
the largest power of $2$ not bigger than $n$, and $\ell=n-2^m$. For example,
$64\leq100<128$ gives $100=2^6+36$. The integers with the same $m$ form the
blocks $2^m\leq n<2^{m+1}$ separated by lines in Table~\ref{tab:results}.

The three results of this section are the tools we need. In the next section
we use them to prove a formula for $J(n)$.

\clearpage
\section{The Formula and Its Generalisation}

In this section we prove the formula suggested by our tests, using the
results of Section~3. We give two proofs, one by induction on $m$ and one by
induction on $n$. Up to now the game always started with person $1$. In
the second part of this section we drop this assumption and find the survivor
when the game starts with any person.

\subsection{The formula \texorpdfstring{$J(n)=2\ell+1$}{J(n) = 2l + 1}}\label{sec:formula}

\subsubsection{What we observe}

In every line of Table~\ref{tab:results}, the column $J(n)$ is equal to the
column $2\ell+1$. In other words, the survivor depends only on the remainder
$\ell$.

\begin{theo}\label{thm:main}
Let $n\geq1$, written as $n=2^m+\ell$ with $0\leq\ell<2^m$. Then
\[
J(n)=2\ell+1 .
\]
\end{theo}

\subsubsection{Proof by induction on \texorpdfstring{$m$}{m}}

For $m\geq0$, let $P(m)$ be the statement
\[
P(m):\quad\text{for every integer $\ell$ with $0\leq\ell<2^m$,}\quad
J(2^m+\ell)=2\ell+1 .
\]
By Observation~\ref{obs:decomp}, every $n\geq1$ is equal to $2^m+\ell$ for
one pair $(m,\ell)$ with $0\leq\ell<2^m$. So if $P(m)$ is true for every
$m\geq0$, the theorem is proved.

\textbf{Base case: $m=0$.} Then $0\leq\ell<2^0=1$, so $\ell=0$, and
\[
J(2^0+0)=J(1)=1=2\times0+1 .
\]
So $P(0)$ is true.

\textbf{Induction step.} Let $m\geq1$, and assume that $P(m-1)$ is true. Let
$\ell$ be an integer with $0\leq\ell<2^m$. We show that
$J(2^m+\ell)=2\ell+1$, in two cases.

\emph{Case 1: $\ell$ is even}, $\ell=2k$. From $0\leq2k<2^m$ we get
$0\leq k<2^{m-1}$, and
\[
2^m+\ell=2^m+2k=2\left(2^{m-1}+k\right).
\]
By Observation~\ref{obs:even}, then by $P(m-1)$ applied to $2^{m-1}+k$,
\[
J(2^m+\ell)=2J\!\left(2^{m-1}+k\right)-1=2(2k+1)-1=4k+1=2\ell+1 .
\]

\emph{Case 2: $\ell$ is odd}, $\ell=2k+1$. From $0\leq2k+1<2^m$ we get
$2k+1\leq2^m-1$, so $k\leq2^{m-1}-1$, that is $0\leq k<2^{m-1}$. Moreover
\[
2^m+\ell=2^m+2k+1=2\left(2^{m-1}+k\right)+1 .
\]
By Observation~\ref{obs:odd}, then by $P(m-1)$ applied to $2^{m-1}+k$,
\[
J(2^m+\ell)=2J\!\left(2^{m-1}+k\right)+1=2(2k+1)+1=4k+3=2\ell+1 .
\]

In both cases $J(2^m+\ell)=2\ell+1$, so $P(m)$ is true.

\textbf{Conclusion.} $P(0)$ is true, and $P(m-1)$ implies $P(m)$ for every
$m\geq1$. By induction, $P(m)$ is true for every $m\geq0$: for every $n\geq1$,
written $n=2^m+\ell$ with $0\leq\ell<2^m$, we have $J(n)=2\ell+1$.
\hfill$\square$

\textbf{Examples.} $100=2^6+36$ gives $J(100)=73$, the value found by halving
in Section~\ref{sub:halving}. With $2026=2^{10}+1002$, the formula gives
$J(2026)=2\times1002+1=2005$ at once.

Since $\ell=n-2^m$, the formula can also be written without $\ell$:
\[
J(n)=2n-2^{m+1}+1,
\]
where $2^m$ is the largest power of $2$ not bigger than $n$. For $n=10$:
$J(10)=20-16+1=5$.

\subsubsection{A second proof, by induction on \texorpdfstring{$n$}{n}}

Relation~\eqref{eq:onestep} of Section~\ref{sec:tool}, which adds one person
at a time, gives a second proof of Theorem~\ref{thm:main}. We prove
$J(n)=2\ell+1$ by induction on $n$. For $n=1=2^0+0$ we have $J(1)=1$.

Let $n\geq2$, and assume the formula for $n-1$: if $n-1=2^{m'}+\ell'$ with
$0\leq\ell'<2^{m'}$, then $J(n-1)=2\ell'+1$. There are two cases.
\begin{itemize}
  \item If $\ell'\leq2^{m'}-2$, then $n=2^{m'}+(\ell'+1)$ with
        $\ell'+1<2^{m'}$, so $\ell=\ell'+1$. Moreover
        $J(n-1)=2\ell'+1\leq2^{m'}+\ell'-1=n-2$, so by \eqref{eq:onestep}
        $J(n)=J(n-1)+2=2\ell'+3=2\ell+1$.
  \item If $\ell'=2^{m'}-1$, then $n=2^{m'+1}$, so $\ell=0$. Here
        $J(n-1)=2^{m'+1}-1=n-1$, the last person, so by \eqref{eq:onestep}
        $J(n)=1=2\times0+1$.
\end{itemize}
In both cases $J(n)=2\ell+1$. \hfill$\square$

The first proof uses the halving of the circle ($n$ becomes about $n/2$), the
second one adds one person at a time ($n-1$ becomes $n$). They lead to the
same formula, which is a good check.

\subsubsection{What the formula explains}

The formula proves the remarks we made on Table~\ref{tab:results}.

\begin{coro}\label{cor:consequences}
Let $n\geq1$, written as $n=2^m+\ell$ with $0\leq\ell<2^m$.
\begin{enumerate}
  \item[(a)] $J(n)$ is odd, and $1\leq J(n)\leq n$.
  \item[(b)] $J(n)=1$ if and only if $n$ is a power of $2$.
  \item[(c)] $J(n)=n$ if and only if $n=2^{m+1}-1$, as for $n=1,3,7,15$.
  \item[(d)] $J(n+1)=J(n)+2$ if $n+1$ is not a power of $2$, and $J(n+1)=1$
        if $n+1$ is a power of $2$.
\end{enumerate}
\end{coro}

\begin{proof}
By Theorem~\ref{thm:main}, $J(n)=2\ell+1$.
(a) The number $2\ell+1$ is odd and at least $1$. Moreover
$2\ell+1\leq2^m+\ell$ is the same as $\ell\leq2^m-1$, which is true.
(b) $2\ell+1=1$ if and only if $\ell=0$, that is $n=2^m$.
(c) $2\ell+1=2^m+\ell$ if and only if $\ell=2^m-1$, that is $n=2^{m+1}-1$.
(d) If $n+1$ is not a power of $2$, then $\ell+1<2^m$ (otherwise
$n+1=2^{m+1}$), so $n+1=2^m+(\ell+1)$ with the same $m$, and
$J(n+1)=2(\ell+1)+1=J(n)+2$. If $n+1$ is a power of $2$, then $J(n+1)=1$ by
(b).
\end{proof}

Point (a) agrees with Step~1 of Section~\ref{sec:even}: every even number
leaves during the first round. Point (d) is relation~\eqref{eq:onestep} again.
The table below summarises these points. On each block $2^m\leq n<2^{m+1}$,
the survivor starts again at $1$ and goes through the odd numbers in
increasing order, up to the last person $2^{m+1}-1$.
\begin{center}
\renewcommand{\arraystretch}{1.2}
\begin{tabular}{|c|c|c|}
\hline
$m$ & block of $n$ & values of $J(n)$, in order\\ \hline
$0$ & $1$ & $1$\\
$1$ & $2,3$ & $1,3$\\
$2$ & $4,\dots,7$ & $1,3,5,7$\\
$3$ & $8,\dots,15$ & $1,3,5,\dots,15$\\
$4$ & $16,\dots,31$ & $1,3,5,\dots,31$\\ \hline
$m$ & $2^m,\dots,2^{m+1}-1$ & $1,3,5,\dots,2^{m+1}-1$\\ \hline
\end{tabular}
\end{center}
A block has $2^m$ values of $n$, and there are exactly $2^m$ odd numbers
between $1$ and $2^{m+1}-1$, so each odd number appears once in each block.

\subsection{When the game starts with person \texorpdfstring{$s$}{s}}\label{sec:general}

So far, person $1$ was always the first person skipped
(Section~\ref{sec:assumption}). In a real game we may not know who starts. So
we ask: \emph{if the game starts with person $s$, who survives?}

\begin{defi}
Let $n\geq1$ and $1\leq s\leq n$. We write $J_s(n)$ for the survivor of the
game on the circle $1,2,\dots,n$ that starts with person $s$: we skip $s$,
eliminate the next person, and go on with the rules of
Definition~\ref{def:game}. In particular $J_1(n)=J(n)$.
\end{defi}

\subsubsection{\texorpdfstring{An example: $n=10$ and $s=4$}{An example: n = 10 and s = 4}}

The rounds are now counted from person $4$.
\begin{itemize}
  \item Round 1 skips $4$ and eliminates $5,7,9,1,3$. The game is back at
        person $4$.
  \item Round 2 skips $4$ and eliminates $6$, skips $8$ and eliminates $10$,
        skips $2$ and eliminates $4$.
  \item Only $8$ and $2$ are left. Round 3 skips $8$ and eliminates $2$.
\end{itemize}
The survivor is $8$, so $J_4(10)=8$ (Figure~\ref{fig:s4}).

\begin{figure}[!ht]
\centering
\setlength{\tabcolsep}{4pt}
\begin{tabular}{cccc}
\begin{tikzpicture}[baseline=(current bounding box.center)]
  \foreach \p in {1,...,10}
    \node[person] (q\p) at ({90-36*(\p-1)}:1.05) {\p};
  \draw[dashed,thick] (q4) circle ({0.5*\nodesize+0.8mm});
\end{tikzpicture}
& \gamecircle{10}{1.05}{}{5,7,9,1,3}{0}
& \gamecircle{10}{1.05}{5,7,9,1,3}{6,10,4}{0}
& \gamecircle{10}{1.05}{5,7,9,1,3,6,10,4}{2}{8}\\[3pt]
\footnotesize start: $s=4$ & \footnotesize Round 1: 5, 7, 9, 1, 3
& \footnotesize Round 2: 6, 10, 4 & \footnotesize Round 3: 2
\end{tabular}
\caption{The game with $10$ people that starts with person $4$ (dashed):
$J_4(10)=8$.}
\label{fig:s4}
\end{figure}

When person $1$ starts, the survivor is $J(10)=5$. Here the game starts three
places later, and the survivor is also three places later: $5+3=8$.

\subsubsection{Counting places from the first person}

The rule \emph{skip one, eliminate the next} only uses the order of the people
around the circle. So we number the \emph{places} from the person who starts:
place $1$ is person $s$, place $2$ is the person after him, and so on. For
$n=10$ and $s=4$:
\begin{center}
\renewcommand{\arraystretch}{1.2}
\begin{tabular}{|c|c|c|c|c|c|c|c|c|c|c|}
\hline
place & 1&2&3&4&5&6&7&8&9&10\\ \hline
person & 4&5&6&7&8&9&10&1&2&3\\ \hline
\end{tabular}
\end{center}
By Lemma~\ref{lem:places}, the survivor is at place $J(10)=5$, that is
person $8$. In general the survivor is at place $J(n)$, and we only need to
know who stands there. Going forward from $s$ we meet $s+1$, $s+2$, \dots, up
to person $n$, and then we come back to person $1$. So the person at place $i$
is $s+i-1$ if $s+i-1\leq n$, and $s+i-1-n$ otherwise. For instance, place $8$
gives $4+8-1=11>10$, so person $11-10=1$.

\subsubsection{The general formula}

Since the survivor is at place $J(n)=2\ell+1$, we obtain the following
result.

\begin{theo}\label{thm:general}
Let $n\geq1$ and $1\leq s\leq n$, and write $n=2^m+\ell$ with $0\leq\ell<2^m$.
If the game starts with person $s$, the survivor is
\[
J_s(n)=
\begin{cases}
s+2\ell & \text{if } s+2\ell\leq n,\\
s+2\ell-n & \text{if } s+2\ell>n.
\end{cases}
\]
In words: the survivor is the person $2\ell$ places after $s$ around the
circle.
\end{theo}

\begin{proof}
The game that starts with $s$ starts from the list
$(a_1,\dots,a_n)=(s,s+1,\dots,n,1,\dots,s-1)$ of the people in the order of
the circle, where $a_i$ is the person at place $i$. By
Lemma~\ref{lem:places}, the survivor is $a_{J(n)}$, the person at place
$J(n)$, and $J(n)=2\ell+1$ by Theorem~\ref{thm:main}. The person at
place $2\ell+1$ is $s+2\ell$ if $s+2\ell\leq n$, and $s+2\ell-n$ otherwise.

It remains to check that one subtraction of $n$ is always enough, that is
$s+2\ell-n\leq n$. Since $\ell<2^m$, we have $2\ell<2^m+\ell=n$, and with
$s\leq n$ we get $s+2\ell<2n$.
\end{proof}

\textbf{Another way to read the formula.} Since $J(n)=2\ell+1$, we have
$s+2\ell=J(n)+(s-1)$. So
\[
J_s(n)=J(n)+(s-1),\quad\text{minus $n$ if this is bigger than $n$.}
\]
Starting the game $s-1$ places later moves the survivor $s-1$ places later.
This is what we saw for $n=10$ and $s=4$: $J_4(10)=5+3=8$.

\textbf{Checks.} For $s=1$ the formula gives $J_1(n)=1+2\ell=J(n)$, as it should. For
$n=10=2^3+2$ we have $\ell=2$, so the survivor is $4$ places after $s$:
\begin{center}
\renewcommand{\arraystretch}{1.2}
\begin{tabular}{|c|c|c|c|c|c|c|c|c|c|c|}
\hline
$s$ & 1&2&3&4&5&6&7&8&9&10\\ \hline
$J_s(10)$ & 5&6&7&8&9&10&1&2&3&4\\ \hline
\end{tabular}
\end{center}
From $s=7$ on, $s+4>10$ and we subtract $10$: for instance $J_8(10)=8+4-10=2$.
We also checked the formula with a computer simulation of the game, for every
$n\leq200$ and every starting person $s$.

\clearpage
\section{Very Large Circles: Base 2}

To use the formula $J(n)=2\ell+1$, we must first write $n=2^m+\ell$, that is,
find the largest power of $2$ not bigger than $n$. For a large $n$ this is not
immediate in base $10$. In base $2$ it is: the decomposition can be read
directly on the digits of $n$, and the formula becomes a simple operation on
these digits. In this whole section, the game starts with person $1$, as in
Section~\ref{sec:formula}: $J(n)$ is the survivor of the usual game.

\subsection{Reading \texorpdfstring{$2^m$}{2\^{}m} and \texorpdfstring{$\ell$}{l} in base 2}\label{sec:bindecomp}

In base $2$, a number is written with the digits $0$ and $1$, and each digit
stands for a power of $2$. For instance
\[
10=(1010)_2=1\cdot2^3+0\cdot2^2+1\cdot2^1+0\cdot2^0=8+2 .
\]
Let $n\geq1$, with $2^m\leq n<2^{m+1}$. Then $n$ has exactly $m+1$ digits in
base $2$, and the first one is a $1$:
\[
n=(1\,b_{m-1}\,b_{m-2}\cdots b_0)_2
 =2^m+\underbrace{b_{m-1}2^{m-1}+\dots+b_12+b_0}_{=\,\ell}.
\]
Since each $b_i$ is $0$ or $1$, the remainder satisfies
$0\leq\ell\leq2^{m-1}+\dots+2+1=2^m-1<2^m$. So this is exactly the
decomposition of Observation~\ref{obs:decomp}: the first digit gives $2^m$,
and the other digits give $\ell$. For $n=10=(1\,010)_2$ we read $2^3=8$ and
$\ell=(010)_2=2$.

\subsection{The rule: move the first digit to the end}

\begin{theo}\label{thm:binary}
If the game starts with person $1$ and $n=(1\,b_{m-1}\cdots b_1b_0)_2$, then
\[
J(n)=(b_{m-1}\cdots b_1b_0\,1)_2 .
\]
In words: to get $J(n)$, we take the first digit of $n$ in base $2$, which is
a $1$, and move it to the end.
\end{theo}

\begin{proof}
By Section~\ref{sec:bindecomp}, $\ell=(b_{m-1}\cdots b_0)_2=\sum_{i=0}^{m-1}b_i2^i$, and
$J(n)=2\ell+1$ by Theorem~\ref{thm:main}. Multiplying by $2$ moves every digit
one place to the left, and adding $1$ writes a $1$ in the last place:
\[
2\ell+1=\sum_{i=0}^{m-1}b_i2^{i+1}+1
=b_{m-1}2^m+b_{m-2}2^{m-1}+\dots+b_02^1+1\cdot2^0
=(b_{m-1}\cdots b_0\,1)_2 .
\]
\end{proof}

Table~\ref{tab:binary} shows the rule on our sixteen tests. When the digit
after the first one is $0$, the result starts with $0$, and this $0$ is simply
dropped: for $n=8=(1000)_2$ we get $(0001)_2=1$.

\begin{table}[!ht]
\centering
\renewcommand{\arraystretch}{1.1}
\begin{tabular}{|c|c|c|c|}
\hline
$n$ & $n$ in base 2 & $J(n)$ in base 2 & $J(n)$\\ \hline
1 & 1 & 1 & 1\\
2 & 10 & 01 & 1\\
3 & 11 & 11 & 3\\
4 & 100 & 001 & 1\\
5 & 101 & 011 & 3\\
6 & 110 & 101 & 5\\
7 & 111 & 111 & 7\\
8 & 1000 & 0001 & 1\\
9 & 1001 & 0011 & 3\\
10 & 1010 & 0101 & 5\\
11 & 1011 & 0111 & 7\\
12 & 1100 & 1001 & 9\\
13 & 1101 & 1011 & 11\\
14 & 1110 & 1101 & 13\\
15 & 1111 & 1111 & 15\\
16 & 10000 & 00001 & 1\\ \hline
\end{tabular}
\caption{The first digit of $n$ moves to the end. The values agree with
Table~\ref{tab:results}.}
\label{tab:binary}
\end{table}

\subsection{What base 2 shows}

Base $2$ gives another view of Corollary~\ref{cor:consequences}.
\begin{itemize}
  \item \textbf{$J(n)$ and $n$ have the same number of digits $1$.} Moving a
        digit, or dropping a $0$ at the front, does not change the number of
        $1$'s. For example $10=(1010)_2$ and $J(10)=5=(101)_2$ both have two.
  \item \textbf{$J(n)=1$ exactly when $n$ is a power of $2$.} The result
        $(b_{m-1}\cdots b_0\,1)_2$ is equal to $1$ only when all the $b_i$ are
        $0$, that is when $n=(10\cdots0)_2=2^m$.
  \item \textbf{$J(n)=n$ exactly when all the digits of $n$ are $1$.} The
        strings $1\,b_{m-1}\cdots b_1b_0$ and $b_{m-1}\cdots b_1b_0\,1$ are equal
        when $b_{m-1}=1$, $b_{m-2}=b_{m-1}$, \dots, $b_0=b_1$, so when all the
        digits are $1$: $n=(11\cdots1)_2=2^{m+1}-1$, as for $n=3$, $7$, $15$.
\end{itemize}

\Needspace{14\baselineskip}
\subsection{Computing \texorpdfstring{$J(n)$}{J(n)} by hand for a large \texorpdfstring{$n$}{n}}

To find $J(n)$ by hand for a large $n$, we follow four steps:
\begin{enumerate}
  \item write $n$ in base $2$, by dividing by $2$ again and again and keeping
        the remainders;
  \item remove the first digit, which is a $1$;
  \item write a $1$ at the end of what is left;
  \item convert the result back to base $10$.
\end{enumerate}

\textbf{Example: $n=1000$.} We divide by $2$ until we reach $0$:
\begin{center}
\renewcommand{\arraystretch}{1.1}
\begin{tabular}{l*{10}{c}}
number & 1000 & 500 & 250 & 125 & 62 & 31 & 15 & 7 & 3 & 1\\
remainder & 0 & 0 & 0 & 1 & 0 & 1 & 1 & 1 & 1 & 1
\end{tabular}
\end{center}
Reading the remainders from right to left gives $1000=(1111101000)_2$. We move
the first $1$ to the end:
\[
(1\,111101000)_2\ \longrightarrow\ (111101000\,1)_2=512+256+128+64+16+1=977 .
\]
So $J(1000)=977$. With the formula: $1000=512+488$, so
$J(1000)=2\times488+1=977$.

\textbf{Example: $n=2026$.} We have $2026=(11111101010)_2$. Moving the first digit to the end gives
\[
(1\,1111101010)_2\ \longrightarrow\ (1111101010\,1)_2=2005,
\]
so $J(2026)=2005$. This agrees with
$2026=2^{10}+1002$ and $2\times1002+1=2005$.

\textbf{Larger values.} Table~\ref{tab:large} gives a few larger values, all obtained by moving the
first digit. We also checked them with the program of Section~6. In the line
$n=10\,000$, two zeros appear at the front after the move and are dropped.

\noindent\begin{minipage}{\textwidth}\centering
\renewcommand{\arraystretch}{1.15}
\footnotesize\setlength{\tabcolsep}{4pt}
\begin{tabular}{|r|l|r|l|}
\hline
$n$ & $n$ in base 2 & $J(n)$ & $J(n)$ in base 2\\ \hline
$1000$ & 1111101000 & $977$ & 1111010001\\
$2026$ & 11111101010 & $2005$ & 11111010101\\
$10\,000$ & 10011100010000 & $3617$ & 111000100001\\
$10^6$ & 11110100001001000000 & $951\,425$ & 11101000010010000001\\
$2^{20}$ & 1 followed by twenty 0 & $1$ & 1\\
$2^{20}-1$ & twenty digits 1 & $1\,048\,575$ & twenty digits 1\\
$2^{20}+1$ & 1, nineteen 0, 1 & $3$ & 11\\
$10^9$ & 111011100110101100101000000000 & $926\,258\,177$ & 110111001101011001010000000001\\ \hline
\end{tabular}
\captionof{table}{Large values of $J(n)$. In each line, $n$ and $J(n)$ have the same
number of digits $1$.}
\label{tab:large}
\end{minipage}
\par\medskip

For $n=10^6$, there are $20$ digits, so $m=19$, $2^{19}=524\,288$ and
$\ell=10^6-524\,288=475\,712$, which gives $J(10^6)=951\,425$. For $n=10^9$,
the game itself needs $999\,999\,999$ eliminations. With base
$2$, we only handle the $30$ digits of $n$. In general, $n$ has about
$\log_2 n$ digits in base $2$, so the work grows like $\log_2 n$ instead of
$n$.

\clearpage
\section{The Python Code}

We wrote two Python functions during the project. The first one, \texttt{J\_0},
plays the game step by step. The second one, \texttt{J\_1}, uses the formula of
Section~4. In this section we explain how they work, we compare their cost,
and we give a faster version of the formula together with the program we used
to check all our results.

\Needspace{23\baselineskip}
\subsection{\texorpdfstring{\texttt{J\_0}: a simulation of the game}{J\_0: a simulation of the game}}

\begin{lstlisting}[language=iPython, caption={The function \texttt{J\_0}, which simulates the game.}]
def J_0(n):
    if n <= 0:
        return "Please enter a positive integer for n"
    else:
        people = [i for i in range(1, n + 1)]
        print(people)
        while len(people) > 1:
            if len(people) % 2 == 0:  # Even case
                for i in range(len(people)//2):
                    h = people.pop(i + 1)
                    print(people, "\t", h, "is eliminated")
            else:  # Odd case
                for i in range(len(people)//2):
                    h = people.pop(i + 1)
                    print(people, "\t", h, "is eliminated")
                people.insert(0, people.pop(len(people) - 1))
                print(people)
        return f"So {people[0]} is the only survivor"
\end{lstlisting}

\textbf{The list.} The circle is a Python list, \texttt{people}, which starts
as $[1,2,\dots,n]$. The person at index $0$ is the next one to be skipped. If
$n\leq0$ the function returns a message, since the game makes no sense.

\textbf{One pass over the list.} Each turn of the \texttt{while} loop goes once
through the list and eliminates every second person. The instruction
\texttt{people.pop(i + 1)} removes the person at index $i+1$. This index is
the right one because the list shrinks while we go: after $i$ eliminations,
$i$ people have left the front part of the list, so the person who was at
index $2i+1$ is now at index $i+1$. For $n=6$:
\[
[1,2,3,4,5,6]\xrightarrow{\ \text{pop}(1)\ }[1,3,4,5,6]
\xrightarrow{\ \text{pop}(2)\ }[1,3,5,6]
\xrightarrow{\ \text{pop}(3)\ }[1,3,5].
\]
With $2k$ people the pass makes $k$ eliminations and ends exactly at the end
of the list, so the next pass starts again with the person at index $0$, who
is skipped.

\textbf{The odd case.} With $2k+1$ people the pass also makes $k$
eliminations, but the last person of the list is skipped at the end of the
pass. The next one to be eliminated is therefore the first person of the list.
The instruction
\begin{center}
\texttt{people.insert(0, people.pop(len(people) - 1))}
\end{center}
moves the last person to the front. Now he is at index $0$ and is skipped, and
the person after him is eliminated, as in the game. For instance
$[1,5,9]$ becomes $[9,1,5]$, and the next pass eliminates $1$.

\textbf{Output for $n=10$.} The function prints the list after each
elimination:
\Needspace{18\baselineskip}
\begin{lstlisting}[numbers=none, caption={Output of \texttt{J\_0(10)}.}]
[1, 2, 3, 4, 5, 6, 7, 8, 9, 10]
[1, 3, 4, 5, 6, 7, 8, 9, 10]   2 is eliminated
[1, 3, 5, 6, 7, 8, 9, 10]      4 is eliminated
[1, 3, 5, 7, 8, 9, 10]         6 is eliminated
[1, 3, 5, 7, 9, 10]            8 is eliminated
[1, 3, 5, 7, 9]                10 is eliminated
[1, 5, 7, 9]                   3 is eliminated
[1, 5, 9]                      7 is eliminated
[9, 1, 5]
[9, 5]                         1 is eliminated
[5, 9]
[5]                            9 is eliminated
So 5 is the only survivor
\end{lstlisting}
The people are eliminated in the order $2,4,6,8,10,3,7,1,9$, which is the order
of the statement. The passes of the program are not exactly our rounds: in a
circle with an odd number of people, we counted the elimination of the first
person in the current round (Section~\ref{sec:conventions}), while the program
does it at the start of the next pass. This changes nothing to the order of
the eliminations or to the survivor.

\Needspace{19\baselineskip}
\subsection{\texorpdfstring{\texttt{J\_1}: the formula with a slow search}{J\_1: the formula with a slow search}}

\begin{lstlisting}[language=iPython, caption={The function \texttt{J\_1}, which uses the formula $J(n)=2\ell+1$.}]
def J_1(n):
    m = 0
    if n == 1:
        return 1
    elif n <= 0:
        return "Please enter a positive integer for n"
    else:
        # Find the greatest m and the corresponding
        # remainder l such that n = 2^m + l
        for j in reversed(range(n)):
            for i in range(n):
                if n == pow(2, i) + j:
                    m, l = i, j
        return 1 + 2*l
\end{lstlisting}

\textbf{Finding $m$ and $\ell$.} The function tries every pair $(i,j)$ with
$0\leq i,j<n$ and keeps the pairs with $n=2^i+j$. The outer loop takes $j$ in
decreasing order, so the last pair kept has the smallest possible $j$. A
smaller remainder means a larger power of $2$, so the last pair kept is the one
where $2^i$ is the largest power of $2$ not bigger than $n$. This is the
decomposition of Observation~\ref{obs:decomp}, and the function returns
$2\ell+1$.

\textbf{Example.} For $n=10$, the pairs found are $(0,9)$, $(1,8)$, $(2,6)$
and $(3,2)$, since $10=1+9=2+8=4+6=8+2$. The last one is $(3,2)$, so
$m=3$, $\ell=2$, and \texttt{J\_1(10)} returns $5$, like \texttt{J\_0(10)}.

\subsection{The cost of \texttt{J\_0} and \texttt{J\_1}}

\textbf{\texttt{J\_0}.} The game has $n-1$ eliminations. Each \texttt{pop}
may move all the people after the removed one, so one elimination can cost up
to $n$ operations, and the whole game about $n^2$. On top of that, the
function prints the whole list after each elimination, which is also of order
$n^2$ characters. The memory used is of order $n$, for the list.

\textbf{\texttt{J\_1}.} The two loops test $n\times n=n^2$ pairs. Moreover,
\texttt{pow(2, i)} is a number with $i+1$ digits in base $2$, so for large $i$
it is not computed in one step. In practice \texttt{J\_1} is even slower than
\texttt{J\_0}: in our tests, for $n=2000$, \texttt{J\_1} took about $3$
seconds, while \texttt{J\_0} (with the printing switched off) took about a
tenth of a second.

So the formula is fast, but this way of using it is not: almost all the work
goes into the search for $m$ and $\ell$. The next subsection fixes this.

\subsection{Faster versions of the formula}

To find the largest power of $2$ not bigger than $n$, we do not need to try
every pair. We start from $p=1$ and double $p$ as long as $2p\leq n$. This
needs only about $\log_2 n$ steps: $30$ steps for $n=10^9$.

\Needspace{29\baselineskip}
\begin{lstlisting}[language=iPython, caption={Faster functions for $J(n)$ and for a game that starts with person $s$.}]
def J_2(n):
    # I look for the largest power of 2 that is not bigger than n,
    # by doubling p as long as I stay below n
    p = 1
    while 2 * p <= n:
        p = 2 * p
    l = n - p
    return 2 * l + 1


def J_binary(n):
    b = bin(n)[2:]               # for example 10 gives '1010'
    return int(b[1:] + b[0], 2)  # I move the first digit to the end


def J_start(n, s):
    # survivor when the game starts with person s (Section 4.2)
    p = 1
    while 2 * p <= n:
        p = 2 * p
    l = n - p
    if s + 2 * l <= n:
        return s + 2 * l
    return s + 2 * l - n
\end{lstlisting}

The function \texttt{J\_2} is the formula $J(n)=2\ell+1$ with this quick
search for $2^m$. The function \texttt{J\_binary} is the rule of
Theorem~\ref{thm:binary}: \texttt{bin(n)} writes $n$ in base $2$ (with the
prefix \texttt{0b}, which we remove), we move the first digit to the end, and
\texttt{int(..., 2)} converts back to base $10$. The function \texttt{J\_start}
is the formula of Theorem~\ref{thm:general}. These three functions give the
answer for $n=10^9$, and even for $n=10^{100}$, at once.

\subsection{Checking all the results}

To check the formulas, we wrote a second simulation, \texttt{play}, which
prints nothing and accepts any starting person. It uses a \texttt{deque}, a
list from which Python can remove the first element quickly.

\Needspace{24\baselineskip}
\begin{lstlisting}[language=iPython, caption={The simulation used to check every formula.}]
from collections import deque


def play(n, s=1):
    # I play the game without printing anything, to check the formulas
    circle = deque(range(1, n + 1))
    circle.rotate(-(s - 1))      # person s is now the first of the list
    while len(circle) > 1:
        circle.rotate(-1)        # the first person is skipped and goes to the back
        circle.popleft()         # the next one is eliminated
    return circle[0]


def check(N):
    for n in range(1, N + 1):
        assert play(n) == J_2(n) == J_binary(n)
        for s in range(1, n + 1):
            assert play(n, s) == J_start(n, s)
    return "all the checks passed for n = 1, ..., " + str(N)
\end{lstlisting}

In \texttt{play}, the first person of the deque is always the next one to be
skipped. \texttt{rotate(-1)} sends him to the back of the circle, and
\texttt{popleft()} eliminates the person who is now first. The function
\texttt{check(200)} compares, for every $n\leq200$, the game with
\texttt{J\_2} and \texttt{J\_binary}, and for every starting person $s$, the
game with \texttt{J\_start}. All the checks passed. We also checked that
\texttt{J\_0} and \texttt{J\_1} give the same survivor as the formula for
every $n\leq300$.

\subsection{Comparison of the four functions}

\begin{table}[!ht]
\centering
\renewcommand{\arraystretch}{1.2}
\begin{tabular}{|l|l|l|l|}
\hline
Function & Method & Number of steps & $n=10^9$\\ \hline
\texttt{J\_0} & plays the game & about $n^2$ & out of reach\\
\texttt{J\_1} & formula, search over all pairs & more than $n^2$ & out of reach\\
\texttt{J\_2} & formula, doubling search & about $\log_2 n$ & at once\\
\texttt{J\_binary} & first digit moved to the end & about $\log_2 n$ & at once\\ \hline
\end{tabular}
\caption{The four functions for $J(n)$.}
\label{tab:code}
\end{table}

\texttt{J\_0} is the one that helped us most at the start: it shows the
circle after each elimination, so we could compare it with our drawings.
\texttt{J\_1} was our first use of the formula. \texttt{J\_2} and
\texttt{J\_binary} show what a formula brings once it is used well: the answer
no longer depends on playing the game.

\clearpage
\section{Conclusion}

We started from a simple game: $n$ people in a circle, every second person
eliminated, and the question of who is left at the end.

By playing the game by hand for $n=1,\dots,16$, we saw that person $1$ survives
when $n$ is a power of $2$, and that the survivor moves two places each time a
person is added. The key remark was that the game only depends on the places
of the people (Lemma~\ref{lem:places}), so that after one turn of the circle
the game is a smaller copy of itself. It gave the two relations
\[
J(2k)=2J(k)-1,\qquad J(2k+1)=2J(k)+1,
\]
and, by induction on $m$, the formula
\[
J(2^m+\ell)=2\ell+1,\qquad 0\leq\ell<2^m ,
\]
which we proved a second time by induction on $n$. Since $2^m$ is the largest
power of $2$ not bigger than $n$, it can also be written
$J(n)=2n-2^{m+1}+1$ with $m=\lfloor\log_2 n\rfloor$.
We then removed the assumption that person $1$ starts: if the game starts with
person $s$, the survivor is the person $2\ell$ places after $s$. In base $2$,
the formula becomes a rule on the digits: the first digit of $n$ moves to the
end. This gives the survivor of a circle of $10^9$ people at once, where the
game would need almost a billion eliminations. Every result was checked by a
computer simulation.

The work followed the method of the course: test small cases, look for a
pattern, state it precisely, and prove it. Each step was needed. The tests
gave the formula, but they could not prove it: sixteen cases say nothing about
the seventeenth. The proof, in turn, came from looking closely at one round of
the game in the tests. The computer was useful at both ends, first to check
our drawings, then to check the formulas on many more cases than we could do
by hand. It also showed us that a formula is only fast if it is used well:
our first program based on the formula was slower than the simulation.

Some questions remain open for us. The first one is to eliminate every third
person instead of every second one, or more generally every $q$-th person. The
proof of Lemma~\ref{lem:places} still works: one step now sends the first
$q-1$ people to the back of the list and eliminates the next one. It gives
\[
J_q(n)=\big((J_q(n-1)+q-1)\bmod n\big)+1,
\]
where $a\bmod n$ is the remainder of $a$ divided by $n$. For $q=2$ this is
relation~\eqref{eq:onestep}. We checked this relation with a simulation for
$q=3,4,5$ and $n<150$. For $q=3$ the survivors for $n=1,\dots,16$ are
$1,2,2,1,4,1,4,7,1,4,7,10,13,2,5,8$: the survivor goes up by $3$ and comes
back near the start from time to time, but we did not find a formula as
simple as $2\ell+1$. The case $q=2$ is special because the circle is exactly
halved at each round. A last question remains: can we predict the whole
order of the eliminations, and not only the survivor?


\clearpage
\section*{Acknowledgements}
\addcontentsline{toc}{section}{Acknowledgements}
We thank Dr Audace A. V. Dossou-Olory for his course, for his tips, and for
the time and energy he gave us during this week. He showed us how to face a
difficult problem with simple and logical steps: start with small cases, look
for a pattern, make a guess, and prove it. We also thank the tutors for their
help, and our classmates for the discussions that helped us move forward.

{\footnotesize
\begin{thebibliography}{99}
\addcontentsline{toc}{section}{References}
\bibitem{course2026} A. A. V. Dossou-Olory, \emph{Mathematical Problem
Solving: Course Introduction}, lecture slides, African Institute for
Mathematical Sciences (AIMS), Senegal, 2026.
\bibitem{projects2026} \emph{Mathematical Problem Solving: Projects},
project sheet, African Institute for Mathematical Sciences (AIMS), Senegal,
October 28, 2026.
\bibitem{concrete} R. L. Graham, D. E. Knuth, O. Patashnik, \emph{Concrete
Mathematics: A Foundation for Computer Science}, 2nd edition, Addison-Wesley,
1994, Chapter 1 (The Josephus Problem).
\end{thebibliography}
}

\end{document}
